% Exposición académica de las dos narrativas conservadas en el cuaderno.
% Procedencia y correspondencia de contenidos: gestion/mapa_narrativa.md.
% H_lazo y rho_lazo no designan la holonomía nonádica Gamma_9.
\section{Compatibilité structurelle de l’action, de la masse et de l’information opérationnelle}
\label{sec:narrativa-estructural}

La composition des équations permet d’identifier une relation plus précise
que l’équivalence entre masse et énergie ou la proportionnalité entre énergie
et fréquence : une même structure de route détermine une fréquence propre,
une relation entre longueurs et une pondération thermique. Le couple angulaire
participe à ces trois lectures parce qu’il contient le coefficient d’action et
détermine les canaux de transport. L’exposé qui suit conserve
conjointement la construction des objets, leurs opérations de lecture et
les conditions de leurs réalisations physiques.

Cette composition permet, en outre, de caractériser l’empreinte laissée par l’ordre
d’une séquence d’opérations. Cette empreinte admet une expression au moyen
du couple angulaire et une application inverse qui restitue le coefficient
d’action. Dans le domaine constitutif, la réponse du vide et le moment
spectral de Catalan sont réunis au moyen d’une même collection fonctionnelle. Le
contenu explicatif de ces relations réside dans les restrictions qu’elles
imposent conjointement, et non dans la simple accumulation d’identités.

La généalogie APP--TRIT--TPK fournit les objets qui sont composés. Les
lectures ultérieures conservent des aspects distincts d’une même structure :
l’échelle, l’orientation, la phase, la mémoire, l’incidence et la réponse. Leur analyse
conjointe établit quelles relations demeurent déterminées lors du changement de
représentation, quelle information une projection élimine et quelles observations
complémentaires permettent de la restituer. Les démonstrations des compositions
et de leurs domaines sont développées dans les sections démonstratives.

\subsection{L’action dans la composition angulaire et dans la structure électronique}
\label{subsec:narrativa-accion-angular}

La section de Planck intervient avant l’évaluation énergétique finale.
Son coefficient participe à la géométrie angulaire utilisée ultérieurement
par la loi des masses. Le lien, avec les unités conservées, est
\begin{equation}
 \hbar_{\mathrm{ret}}
   =\eta_{\mathrm{ret}}\,10^{-34}\mathcal U_S,
 \qquad
 A_{\deg}=1000\alpha,
 \qquad
 C^*_{\deg}=2(\eta_{\mathrm{ret}}+\alpha).
 \label{eq:narrativa-par-angular}
\end{equation}
Le coefficient $\eta_{\mathrm{ret}}$ est adimensionnel et peut donc
être composé avec $\alpha$ dans l’expression angulaire. La constante dimensionnelle
$\hbar_{\mathrm{ret}}$ conserve son unité d’action. Cette distinction permet de
suivre une même dépendance sans additionner des grandeurs de dimensions différentes.

L’action remplit ainsi deux fonctions liées. D’une part, elle fixe une
échelle dimensionnelle ; d’autre part, elle participe à la structure qui détermine la
répartition du transport entre les deux orientations. Après la conversion
unique des degrés en radians, l’angle et le contre-angle produisent les canaux
$q_+$ et $q_-$. Ces canaux interviennent dans les incidences hexadiques--octadiques,
dans la fonctionnelle de Barbero--Immirzi, dans la réponse constitutive et dans les
termes de la loi des masses.

La substitution de son expression structurelle à une constante rend visibles
des dépendances qui demeurent implicites lorsque l’équation est écrite au moyen
de constantes isolées. L’expression développée transmet l’opération qui
détermine la valeur, les coordonnées dont elle dépend et les conditions
dans lesquelles elle peut être réutilisée. Sa fonction explicative procède de cette
information transmise à l’étape suivante.

Dans l’électron, l’expression complète conserve l’exposant
$\varphi/\pi^2$, les corrections construites avec $\alpha$ et $\Delta_4$,
et le facteur de retour
\begin{equation}
 \left(\frac{H_5}{\eta_{\mathrm{ret}}}\right)^{80-54\alpha+6\Delta_4}.
 \label{eq:narrativa-factor-electronico}
\end{equation}
Le même $\eta_{\mathrm{ret}}$ qui intervient dans le contre-angle participe
à la lecture électronique complète. La géométrie angulaire et la correction
d’action de l’électron partagent donc une dépendance effective.
Étudier conjointement les deux équations permet de déterminer quelles relations
sont imposées par l’intervention d’une même coordonnée structurelle
dans des opérations distinctes.

\subsection{Fréquence propre et normalisation structurelle de la masse}
\label{subsec:narrativa-masa-frecuencia}

La loi générale des masses conserve son préfacteur électronique. On désigne par
$\Xi_\Gamma$ le facteur structurel complet d’une route $\Gamma$ : il réunit la
normalisation électronique, le caractère de la route avec sa signature et ses tours,
et la puissance correspondante du rapport d’action
$H_5/\eta_{\mathrm{ret}}$. Cette notation abrège une composition
d’opérations déjà définies ; son contenu reste lié à ces opérations.

L’horloge commune a pour fréquence angulaire
$\omega_0=2\pi/(108t_0)$, avec $t_0>0$. La composition de la loi des masses avec la section
d’action donne
\begin{equation}
 \boxed{
 m_\Gamma=\frac{\hbar_{\mathrm{ret}}\omega_0}{c^2}\,\Xi_\Gamma,
 \qquad
 E_\Gamma=\hbar_{\mathrm{ret}}\omega_0\Xi_\Gamma,
 \qquad
 \Omega_\Gamma=\frac{E_\Gamma}{\hbar_{\mathrm{ret}}}
             =\omega_0\Xi_\Gamma.}
 \label{eq:narrativa-masa-energia-frecuencia}
\end{equation}
Ici, $E_\Gamma$ est l’énergie au repos et $\Omega_\Gamma$ sa fréquence
énergétique associée.

Cette lecture sépare le rythme commun de la contribution propre à chaque
structure. Deux sections peuvent partager la période élémentaire et présenter
des masses et des fréquences propres différentes : leurs signatures, retours et mémoires
produisent des facteurs $\Xi_\Gamma$ différents. La diversité des lectures
spécifiques se trouve rapportée à la structure de la section, tandis que l’horloge
commune établit leur normalisation temporelle.

Dans le quotient $E_\Gamma/\hbar_{\mathrm{ret}}$, le facteur
absolu d’action qui apparaît explicitement dans
\eqref{eq:narrativa-masa-energia-frecuencia} se simplifie. Toutefois, sa
participation structurelle aux canaux et au rapport de retour contenu
dans $\Xi_\Gamma$ demeure. La simplification d’un facteur dans une lecture ultérieure
doit être distinguée de l’élimination des relations au moyen desquelles
ce facteur a été construit et est intervenu de nouveau.

Les rapports entre masses coïncident avec les rapports entre ces fréquences
propres. Une modification commune des unités d’action, de temps ou de vitesse
laisse ces rapports inchangés. La structure relative a donc
une exigence de confrontation propre : ses relations entre sections doivent
rester cohérentes, indépendamment de la normalisation commune choisie.

La fréquence associée à une phase relevée conserve, en outre, le nombre de
tours de l’histoire. Deux phases égales après leur réduction modulo
$2\pi$ peuvent correspondre à des relèvements différents. Le retour de phase
conserve une régularité, tandis que l’histoire distingue le parcours qui
a produit ce retour. L’identification de l’état complet à la phase
réduite éliminerait précisément cette distinction.

\subsection{Dilatation réciproque des lectures spatiales de la masse}
\label{subsec:narrativa-reciprocidad-espacial}

Dans la réalisation gravitationnelle circulaire, l’action et la gravitation
satisfont
\begin{equation}
 G=\frac{c^3L_\alpha^2}{\hbar_{\mathrm{ret}}},
 \qquad R_{108}=L_\alpha=\frac{c}{\omega_0}.
 \label{eq:narrativa-gravedad-accion}
\end{equation}
La composition avec la masse de
\eqref{eq:narrativa-masa-energia-frecuencia} détermine
\begin{equation}
 \boxed{
 r_{g,\Gamma}=R_{108}\Xi_\Gamma,
 \qquad
 \bar\lambda_\Gamma=\frac{R_{108}}{\Xi_\Gamma},
 \qquad
 r_{g,\Gamma}\bar\lambda_\Gamma=R_{108}^{2}.}
 \label{eq:narrativa-longitudes-reciprocas}
\end{equation}
On utilise la convention du rayon gravitationnel réduit du développement ;
la construction de $L_\alpha$ et la normalisation conjointe sont
démontrées dans la section~\ref{md:rigidez}. L’horloge et le rayon sont
deux expressions compatibles de cette même longueur.
$\bar\lambda_\Gamma$ désigne la longueur de Compton réduite. Ces
identités réciproques sont formulées pour $\Xi_\Gamma>0$ ; la lecture
$R_{108}/\Xi_\Gamma$ requiert ce domaine massique positif.

Le même facteur qui augmente la fréquence propre dilate la lecture
gravitationnelle et contracte la longueur associée à l’action et à la masse. Les deux
transformations conservent une aire commune. Dans cette réalisation, le facteur
massique agit comme paramètre de dilatation réciproque : une longueur croît avec
lui et l’autre diminue inversement.

Cette relation fournit un contenu mathématique précis au lien
entre masse et configuration spatiale. La hiérarchie de la composition est
essentielle : on détermine d’abord la route et son facteur structurel ; on obtient
ensuite ses lectures temporelles, énergétiques et spatiales. La structure
est transmise aux lectures au moyen d’opérations déterminées, au lieu
d’être reconstruite rétrospectivement à partir d’une masse mesurée.

La condition circulaire appartient au résultat
\eqref{eq:narrativa-longitudes-reciprocas}. Les deux longueurs sont des lectures
définies par cette réalisation ; leur interprétation comme rayons matériels
d’une particule requerrait l’identification physique correspondante.
L’égalité établit ici une réciprocité exacte entre les grandeurs
définies, avec le système de coordonnées d’action, l’horloge et la vitesse conservés.

L’incorporation de la quantité de mouvement maintient cette masse structurelle dans l’équation
de propagation. Dans la réalisation libre correspondante, la fréquence
propre détermine le terme massique de la dispersion, tandis que la quantité de mouvement
décrit la variation spatiale. Ainsi sont différenciés la structure
persistante de la section et son état de mouvement.

\subsection{Orientation des canaux et réponse constitutive du vide}
\label{subsec:narrativa-vacio-orientacion}

La composition de la permittivité, de la perméabilité et de la fonctionnelle de Barbero--Immirzi
permet de distinguer la valeur d’une réponse de la structure qui la
détermine. Le lecteur constitutif utilise deux canaux étiquetés, de
réponses $r_+$ et $r_-$. La vitesse normalisée dépend de leur produit ;
l’impédance dépend de leur quotient. La vitesse conserve une relation
commune entre les canaux, tandis que l’impédance distingue la répartition de la
réponse entre leurs orientations.

Au sein de la collection normalisée, la vitesse et la fonctionnelle de
Barbero--Immirzi permettent de retrouver conjointement les deux réponses
étiquetées. La fonctionnelle apporte une information orientée que la vitesse
isolée ne détermine pas. La restitution s’appuie sur la composition complète
de leurs lecteurs et conserve la distinction entre chaque canal et leur agrégat.

Soient $x$ et $y$ les coordonnées en radians de l’angle et du contre-angle.
Pour $x>0$ et $|y|<x$, la réponse satisfait
\begin{equation}
 \widehat c(x,y)\geq\widehat c(x,0),
 \label{eq:narrativa-convexidad-velocidad}
\end{equation}
avec égalité uniquement lorsque $y=0$. Plus précisément, en maintenant $x$
fixe, $\log\widehat c$ est une fonction paire et strictement convexe de $y$.
La démonstration de cette propriété appartient au développement constitutif
ultérieur.

La vitesse perd le signe de l’orientation, mais conserve un effet
de la séparation entre les canaux. Près de la configuration symétrique,
cet effet apparaît au second ordre ; l’impédance distingue l’orientation
dès le premier ordre. Par conséquent, remplacer les deux canaux par leur moyenne avant
d’évaluer la réponse change le résultat. Une lecture scalaire peut masquer
une orientation et conserver, en même temps, une conséquence de celle-ci.

Ce comportement spécifie quelle information chaque opération transmet.
La perte du signe, la conservation d’une séparation et la récupération
de la répartition constituent des propriétés différentes. La structure maintient
des conséquences observables même lorsque l’une de ses lectures est exprimée
au moyen d’un seul nombre.

L’inégalité \eqref{eq:narrativa-convexidad-velocidad} décrit la collection
mathématique de réponse dans son domaine. Une évolution physique doit se situer
dans les états et les trajectoires effectivement produits par TPK. La
paramétrisation analytique permet d’étudier la sensibilité du lecteur ;
l’identification d’une variation paramétrique à une trajectoire physique
conserve son obligation de réalisation.

\subsection{Boltzmann, fréquence structurelle et pondération thermique ternaire}
\label{subsec:narrativa-boltzmann}

La section thermique relie l’énergie du cycle à la constante de
Boltzmann et à sa température correspondante :
\begin{equation}
 k_B\Theta_{108}\log3=E_{\mathrm{clk}}.
 \label{eq:narrativa-boltzmann-ciclo}
\end{equation}
Comme l’énergie d’une section est
$E_\Gamma=E_{\mathrm{clk}}\Xi_\Gamma$, le poids thermique évalué à cette
température est
\begin{equation}
 \boxed{\exp\!\left(-\frac{E_\Gamma}{k_B\Theta_{108}}\right)
       =3^{-\Xi_\Gamma}.}
 \label{eq:narrativa-peso-termico}
\end{equation}
La base ternaire apparaît par composition de la normalisation informationnelle
avec l’action et l’horloge. La structure de route qui détermine la masse
détermine également l’exposant du poids thermique.

Le poids et la probabilité remplissent des fonctions distinctes. Les probabilités
s’obtiennent en normalisant les poids et en conservant les multiplicités des
états. De même, plusieurs histoires équivalentes requièrent une identification
d’états avant de recevoir une multiplicité physique : le nombre d’histoires
et la dégénérescence de l’état correspondent à des objets différents.

La composition permet de réunir les lectures dans une identité commune :
\begin{equation}
 \boxed{
 \Xi_\Gamma
 =\frac{\Omega_\Gamma}{\omega_0}
 =\frac{r_{g,\Gamma}}{R_{108}}
 =\frac{R_{108}}{\bar\lambda_\Gamma}
 =\frac{E_\Gamma}{k_B\Theta_{108}\log3}.}
 \label{eq:narrativa-cuatro-lecturas}
\end{equation}
On obtient quatre manières coordonnées de lire une même composante
structurelle : temporelle, spatiale, gravitationnelle et thermique. Chaque égalité
conserve les conditions de sa réalisation, en particulier la condition
circulaire gravitationnelle et la température du cycle.

L’identité \eqref{eq:narrativa-cuatro-lecturas} spécifie la relation
que partagent les grandeurs et fournit un critère pour confronter leur
compatibilité. La masse relative peut être caractérisée au moyen de la fréquence
relative ; la même coordonnée détermine une dilatation réciproque des
longueurs et une pondération thermique. Les fonctions de lecture sont distinctes,
mais leur dépendance se trouve liée par la généalogie commune.

La relation thermique identifie l’état et la normalisation dans lesquels
intervient $\Theta_{108}$. Cette identification maintient séparées l’énergie
de la section et son éventuel transfert sous forme de chaleur ; elle maintient également
localisée la signification de la température du cycle, sans l’étendre à
l’univers comme température unique. La capacité explicative procède des
conditions explicites de la composition.

\subsection{Ellipse d’action et transport de la forme géométrique}
\label{subsec:narrativa-elipse-transporte}

L’ellipse d’action a des demi-axes dont le produit est $2\hbar$. L’aire
orientée accumulée pendant un tour est $h=2\pi\hbar$. Une déformation
elliptique redistribue les demi-axes en conservant leur produit. Dans la
réalisation électrique et magnétique, on obtient deux contributions
proportionnelles à $\cos^2\theta$ et $\sin^2\theta$, dont la somme est
$\hbar\omega$. La structure détermine donc la répartition de
l’énergie pendant le cycle, en plus de sa valeur totale.

Le transport d’une ellipse dont la forme change exige de conserver la relation
entre la déformation et l’action. Soit $\zeta$ la coordonnée de déformation,
déterminée par le rapport entre le contre-angle et l’angle. Pour une
déformation différentiable, le transport exact introduit
\begin{equation}
 H_{\mathrm{tot}}
 =\frac{\omega}{2}
   \left(e^{-\zeta}Q^2+e^\zeta P^2\right)
  +\frac{\dot\zeta}{2}QP.
 \label{eq:narrativa-hamiltoniano-transporte}
\end{equation}
Le premier terme décrit l’évolution sur une ellipse de forme fixe.
Le second procède du transport d’une forme variable. La démonstration
de l’identité de transport établit que ce terme compense
exactement la variation géométrique et conserve l’action quadratique
correspondante. Son omission ou l’inversion de son signe détruisent cette
compensation.

L’interprétation exige de distinguer deux situations. Lorsqu’on change
uniquement de coordonnées, le terme supplémentaire est une connexion de
représentation. Un échange physique requiert, en revanche, une variation
relative entre système et référence. La réexpression d’un Hamiltonien
conserve le contenu de la représentation ; l’attribution d’une interaction
nécessite de déterminer les opérations relatives qui agissent sur le système.

Le couple angulaire directeur conserve une autre dépendance décisive : fixer $\alpha$
et la section d’action fixe également $\zeta$. Une évolution de la forme
doit donc être extraite des états et des opérations effectifs. La
déformation paramétrique à elle seule décrit une collection de représentations ;
son identification comme évolution HMT transmet les relations qui produisent
l’angle et le contre-angle.

Le transport conjoint de la forme et de la section d’action conserve
l’action normalisée, tandis que l’aire non normalisée change selon
le quotient entre les deux sections. Cette construction distingue trois
opérations : redistribuer une action conservée, transporter une section
d’action et changer la représentation ou les unités. La distinction permet
d’examiner les interactions en conservant un bilan qui inclut l’origine de la
variation, au lieu d’attribuer une génération d’énergie à une simple déformation
de coordonnées.

\subsection{Compatibilité conjointe et confrontation relationnelle}
\label{subsec:narrativa-compatibilidad-conjunta}

L’axe de recherche est la détermination des conditions de
compatibilité qu’une même structure impose à ses différentes
manifestations physiques. Deux lignes de composition se complètent.

La première réunit des relations statiques : les observables complémentaires
restituent les canaux, et la structure retrouvée détermine l’action,
la réponse constitutive et les facteurs de masse. La confrontation pertinente
est conjointe. Une modification qui préserve un chiffre doit également conserver
les autres relations qu’utilise la même construction. La vitesse,
Barbero--Immirzi et l’impédance fournissent un système de vérification ;
les rapports de masses et de fréquences en fournissent un autre.

La seconde ligne étudie les histoires effectives et compare leurs résultats
après avoir transporté de manière cohérente le système et sa référence. Les
différences de phase relative, de réponse constitutive relative ou de transfert
entre modes ont un contenu distinct des différences dues au
nom des coordonnées. La mémoire, l’orientation et les retours
interviennent au moyen des opérations qu’une lecture scalaire peut éliminer.

Une relation structurelle acquiert une capacité explicative et de confrontation
lorsque plusieurs lectures indépendantes doivent concorder parce qu’elles procèdent d’une
même structure. Sa valeur réside dans la rigidité de cette concordance :
l’action, le rythme, la réponse spatiale et la pondération informationnelle cessent de
pouvoir être choisis séparément au sein de la réalisation déclarée. La
composition réunit ces restrictions et permet de démontrer leurs conséquences
sans remplacer leur généalogie par des ajustements individuels des résultats.

\subsection{Clôture de phase et holonomie relative de l’état}
\label{subsec:narrativa-holonomia-relativa}

Le cycle dodécaphasique contient le plan dans lequel est réalisée la racine
interne de $-1$. L’inclusion de ce plan dans les douze composantes est
explicite ; le quart de tour utilisé ci-dessous procède de cette
construction. Sur celui-ci agit la forme de l’ellipse d’action déterminée
par le couple angulaire, dont la coordonnée est
\begin{equation}
 \xi=\frac{C^*}{A},\qquad
 \zeta=\operatorname{artanh}\xi.
 \label{eq:narrativa-coordenada-forma}
\end{equation}
Le quotient utilise les deux angles exprimés dans la même unité et a
la même valeur en degrés et en radians. La coordonnée réelle de forme est
définie pour $|\xi|<1$. La forme conjuguée correspond au
changement de signe de $\zeta$.

On compose le quart de tour d’une forme avec celui de sa conjuguée et avec
leurs inverses. Si $U_+$ et $U_-$ désignent les deux opérations, la
composition exacte est
\begin{equation}
 \boxed{
 H_{\mathrm{lazo}}=U_+U_-U_+^{-1}U_-^{-1}
 =\begin{pmatrix}\rho_{\mathrm{lazo}}&0\\[2pt]
                  0&\rho_{\mathrm{lazo}}^{-1}\end{pmatrix},
 \qquad
 \rho_{\mathrm{lazo}}=
 \left(\frac{1+\xi}{1-\xi}\right)^2.}
 \label{eq:narrativa-holonomia-lazo}
\end{equation}
La composition dilate une coordonnée et contracte la conjuguée par le facteur
inverse : elle conserve l’aire et modifie l’état. Pour $\xi\ne0$, son
résultat diffère de l’identité.

Les quatre opérations comprennent leurs inverses et leurs accroissements de phase
déclarés ont une somme nulle. L’ordre de composition conserve toutefois un
effet. Une description qui ne retient que les accroissements qui s’ajoutent et
se compensent élimine cet effet relatif. La clôture de la comptabilité
de phase et le retour de l’état complet sont donc des propriétés
distinctes de la séquence.

L’égalité admet un relèvement à l’espace initial à douze
composantes. L’opérateur agit par la dilatation réciproque sur le
plan considéré et comme l’identité sur son complément. Le lien avec
le cycle dodécaphasique est réalisé au moyen d’une transformation explicite ;
le plan conserve ainsi sa provenance au sein de l’espace des phases.

La notation $H_{\mathrm{lazo}}$ désigne cette composition ultérieure et
maintient distincte l’holonomie nonadique $\Gamma_9$. La réalisation
illustre de manière concrète une distinction structurelle : la clôture d’une
lecture du parcours peut coexister avec une transformation persistante
de l’état. La provenance des opérateurs et la portée de leur
relèvement demeurent exposées dans le développement démonstratif.

\subsection{Restitution du coefficient d’action à partir du transport relatif}
\label{subsec:narrativa-restitucion-accion}

Le couple angulaire de \eqref{eq:narrativa-par-angular} peut être exprimé avec
$s_0=10^{-34}\mathcal U_S$ et
$\hbar_{\mathrm{ret}}=s_0\eta_{\mathrm{ret}}$. Le coefficient
$\eta_{\mathrm{ret}}$ conserve son caractère adimensionnel, tandis que $s_0$
conserve l’échelle dimensionnelle. En incorporant le contre-angle complet
dans \eqref{eq:narrativa-holonomia-lazo}, on obtient
\begin{equation}
 \boxed{
 \rho_{\mathrm{lazo}}=
 \left(\frac{501\alpha+\eta_{\mathrm{ret}}}
              {499\alpha-\eta_{\mathrm{ret}}}\right)^2.}
 \label{eq:narrativa-lazo-alfa-accion}
\end{equation}
Cette expression utilise la branche positive
$0<\eta_{\mathrm{ret}}<499\alpha$. Les nombres $501$ et $499$ procèdent
de $500\pm1$, et $500$ procède de $1000/2$ : ce sont des conséquences des
normalisations qui interviennent déjà dans le couple angulaire.

La relation entre $\alpha$ et le coefficient d’action détermine la
séparation entre les réponses conjuguées après le parcours
composé. L’action participe aussi bien à l’échelle dimensionnelle qu’à
la forme qui gouverne le transport. La formule rend explicite cette
double intervention sans identifier les deux fonctions.

La relation admet l’inversion
\begin{equation}
 \boxed{
 \eta_{\mathrm{ret}}=
 \alpha\left[
 500\,\frac{\sqrt{\rho_{\mathrm{lazo}}}-1}
               {\sqrt{\rho_{\mathrm{lazo}}}+1}-1\right].}
 \label{eq:narrativa-inversa-accion}
\end{equation}
Avec $\alpha$ et les axes orientés conservés, la réponse de la boucle
permet de restituer le coefficient d’action. La lecture est strictement
croissante par rapport à $\rho_{\mathrm{lazo}}$ sur la branche déclarée.

On obtient une caractérisation opérationnelle : le coefficient d’action
peut être restitué au moyen de la dilatation relative produite par les
formes conjuguées du couple angulaire. Sa définition génératrice demeure
dans la construction précédente ; la caractérisation opérationnelle établit
comment le retrouver dans une représentation distincte.

L’échelle $s_0$ reste nécessaire pour retrouver la constante
dimensionnelle $\hbar_{\mathrm{ret}}$. La forme du transport conserve
une relation adimensionnelle ; l’échelle d’action détermine la
normalisation physique complète. Cette séparation identifie avec précision
quelle donnée l’inversion restitue et quelle structure supplémentaire transmet la
lecture dimensionnelle.

\subsection{Itération de l’holonomie et information opérationnelle de l’histoire}
\label{subsec:narrativa-informacion-operacional}

La répétition de la boucle conserve le même programme d’opérations et
produit une loi explicite d’accumulation :
\begin{equation}
 H_{\mathrm{lazo}}^n=
 \begin{pmatrix}\rho_{\mathrm{lazo}}^n&0\\
                 0&\rho_{\mathrm{lazo}}^{-n}\end{pmatrix}.
 \label{eq:narrativa-iteracion-lazo}
\end{equation}
Une préparation $(Q_0,P_0)$ se transforme en
\begin{equation}
 Q_n=\rho_{\mathrm{lazo}}^nQ_0,
 \qquad
 P_n=\rho_{\mathrm{lazo}}^{-n}P_0.
 \label{eq:narrativa-iteracion-estado}
\end{equation}
Le produit $Q_nP_n$ demeure constant et le transport conserve
l’aire symplectique, tandis que la proportion entre les deux
coordonnées change.

Une règle finie répétée peut laisser une empreinte cumulative qui
distingue ses itérations. Une préparation connue et une référence
conservée permettent de retrouver $n$ au moyen de l’accroissement logarithmique
d’une coordonnée non nulle, pour $\rho_{\mathrm{lazo}}\ne1$.
Lorsque seul l’état final est
conservé et que la préparation est inconnue, la même
restitution reste indéterminée.

L’information accessible dépend de la lecture retenue : bilan,
amplitude, proportion orientée, préparation ou opérateur complet. Chaque
choix peut identifier des histoires qu’une autre lecture distingue.
La comparaison de
\begin{equation}
 U_+U_-U_+^{-1}U_-^{-1}
 \qquad\text{et}\qquad
 U_+U_+^{-1}U_-U_-^{-1}
 \label{eq:narrativa-orden-multiplicidades}
\end{equation}
fournit une confrontation directe. Les deux produits contiennent les mêmes
opérations avec les mêmes multiplicités. Le second donne l’identité ;
le premier donne $H_{\mathrm{lazo}}$. Une statistique des fréquences
d’apparition est insuffisante pour prédire cette réponse : il faut
également conserver l’ordre.

Ce résultat précise quelle description statistique est insuffisante
pour l’objet étudié et quelles observations restituent la différence.
La théorie de l’information de Shannon admet des distributions sur des
séquences ordonnées et peut représenter cet ordre. La distinction
opérationnelle se réfère au contenu retenu par l’observable choisi.

L’information opérationnelle d’une histoire est caractérisée par
sa capacité à produire des réponses distinguables pour des préparations
et des détecteurs spécifiés. Cette définition permet de préciser combien
d’information une lecture conserve. Sa portée correspond à la représentation
utilisée : des histoires distinctes peuvent induire un même opérateur
réduit et conserver des différences dans d’autres composantes de la généalogie
TPK.

La trace de $H_{\mathrm{lazo}}$ coïncide avec celle de
$H_{\mathrm{lazo}}^{-1}$. Elle détecte l’intensité de l’effet, mais élimine
son orientation. Une lecture qui conserve les axes et compare leurs
réponses distingue les deux sens. Le choix du détecteur fait donc
partie de la restitution du parcours.

La finitude du programme, la croissance d’une mémoire observable et
un taux entropique décrivent des propriétés différentes. Un taux de
croissance combinatoire nul peut coexister avec des états distinguables.
L’entropie de la dynamique complète requiert de spécifier son espace
d’états et, le cas échéant, sa mesure. La composition démontre
l’accumulation et sa restitution dans la représentation déclarée ;
l’évaluation d’une entropie thermodynamique conserve ses variables
et ses conditions physiques propres.

\subsection{Potentiel spectral et réciprocité différentielle des réponses}
\label{subsec:narrativa-potencial-espectral}

Les canaux constitutifs conservent les coordonnées
$x=A_{\mathrm{rad}}$, $y=C^*_{\mathrm{rad}}$ et les incidences $90/120$.
Leur composition détermine l’identité
\begin{equation}
 \boxed{
 \frac{\partial\log\widehat c}{\partial y}
 =\frac{\partial\log\widehat Z}{\partial x}.}
 \label{eq:narrativa-reciprocidad-diferencial}
\end{equation}
Les grandeurs $\widehat c$ et $\widehat Z$ désignent la vitesse et
l’impédance normalisées du lecteur.

La sensibilité de la propagation à la séparation orientée se trouve
liée à la sensibilité de l’impédance à la coordonnée commune.
Les deux réponses partagent une fonction spectrale et leurs variations
doivent maintenir cette relation pour continuer d’appartenir à la même
construction.

Il existe une fonction $\mathscr F(x,y)$ dont les dérivées sont
\begin{equation}
 \partial_x\mathscr F=\log\widehat c,
 \qquad
 \partial_y\mathscr F=\log\widehat Z.
 \label{eq:narrativa-gradiente-potencial}
\end{equation}
Il s’agit d’un potentiel spectral défini à partir des canaux $90/120$,
avec une série convergente et une normalisation explicite. Son hessien est
défini négatif dans $x>|y|$, propriété qui permet de démontrer
l’unicité et la sensibilité de la restitution.

L’égalité entre les dérivées croisées établit une restriction
sur la variation conjointe des réponses au sein de cette collection
de réalisation. Une modification qui détruit l’identité
\eqref{eq:narrativa-reciprocidad-diferencial} change le lecteur
constitutif. La confrontation porte ainsi sur les relations de variation,
en plus des valeurs évaluées.

Le potentiel est univoque : l’intégrale de sa différentielle sur
une boucle fermée dans les deux coordonnées s’annule. La composition
ordonnée des transports peut conserver, en revanche, une holonomie
non triviale. La projection constitutive et le transport opératoriel
retiennent des aspects différents de l’histoire ; leur comparaison exige de
conserver leurs domaines et leurs opérations de lecture.

\subsection{Moment de Catalan et réponse contractive des incidences}
\label{subsec:narrativa-catalan-potencial}

La connexion entre le moment orienté de Catalan et le potentiel
constitutif est exprimée au moyen du moment de profondeur deux
\begin{equation}
 D_2(z)=\sum_{n\geq1}\frac{z^n}{n^2}.
 \label{eq:narrativa-momento-dos}
\end{equation}
Le moment de Catalan s’obtient comme partie imaginaire de
$D_2(is)$ ; en $s=1$, il retrouve la constante de Catalan. Le potentiel
constitutif utilise la même fonction sur les contractions des
canaux, en $q_\pm^{90}$ et $q_\pm^{120}$, avec les normalisations
correspondantes.

Le quart de tour orienté et la réponse contractive $90/120$
sont réunis au moyen d’un même moment spectral, évalué sur des
supports distincts. Catalan acquiert ici sa signification par
l’opération spectrale qu’il représente et par la relation de cette
opération avec la réponse constitutive. La collection fonctionnelle
conserve une information qui dépasse celle de sa valeur extrême.

La composition maintient également les différences entre ses objets.
La valeur isolée de Catalan ne détermine pas tous les canaux ;
la relation utilise le moment spectral complet, ses arguments
et l’orientation. La convergence de la série, les identités
de ses dérivées et sa compatibilité avec l’amplification font
partie du développement démonstratif de cette connexion.

\subsection{Restitution exacte et stabilité face à une précision finie}
\label{subsec:narrativa-estabilidad-restitucion}

Le couple $(\widehat c,\widehat Z)$ retrouve de manière unique les deux
paramètres du lecteur dans le domaine déclaré. Les bornes de
sensibilité relient, en outre, l’erreur des réponses à
l’erreur des paramètres restitués. L’unicité exacte et la
stabilité de l’inversion constituent des résultats distincts et
complémentaires.

Cette distinction sépare une information éliminée par la projection
d’une information conservée mais difficile à résoudre avec une précision
finie. Dans les régions de contractions intenses, des variations appréciables
des paramètres peuvent produire des différences très petites dans les
réponses. L’inversion reste unique pour des données exactes,
tandis qu’une incertitude finie peut empêcher de distinguer ces
variations.

La restitution généalogique exige de spécifier aussi bien quelle structure
peut être retrouvée que la stabilité avec laquelle elle peut l’être.
La marge de restitution désigne la borne qui quantifie cette stabilité
dans un domaine. Le choix de l’observable doit tenir compte de la différence
que l’on cherche à résoudre : une lecture peut être correcte et s’avérer
peu sensible à celle-ci.

La trace de la boucle offre un exemple de cette sélection. Sa lecture
élimine le sens du parcours ; la réponse sur des axes orientés
le conserve. La détermination du détecteur adéquat appartient à
la démonstration de récupérabilité et à la conception de la confrontation.

\subsection{Synthèse : coordonnées de la section et mémoire des opérations}
\label{subsec:narrativa-sintesis}

L’identité \eqref{eq:narrativa-cuatro-lecturas} conserve comme axe
la relation entre le facteur structurel d’une section, sa fréquence
propre, les deux longueurs réciproques et sa lecture thermique. Les conditions
demeurent explicites : le même système de coordonnées d’action, la réalisation
gravitationnelle circulaire et la température du cycle.

La composition opératorielle ajoute une distinction centrale. Le facteur
de la section détermine ses lectures coordonnées ; le transport
détermine le résultat de la composition d’opérations sur celle-ci et
l’empreinte que conserve leur ordre. Le facteur massique et la dilatation
$\rho_{\mathrm{lazo}}$ sont des objets différents. Leur relation doit
passer par les lecteurs correspondants pour distinguer une
transformation de l’état d’une identification massique.

Planck intervient dans la normalisation de l’action et dans la forme
angulaire du transport. Boltzmann intervient dans la relation entre
l’énergie et la pondération thermique. La réciprocité gravitationnelle conserve
le produit des longueurs dans sa réalisation déclarée. La vitesse
et l’impédance fournissent des lectures complémentaires des canaux.
Cette architecture admet conjointement une description des
compatibilités et une description de la mémoire opérationnelle.

La capacité explicative de la structure réside dans la détermination de
plusieurs réponses, l’imposition de relations entre elles et la spécification de
l’histoire qui peut être restituée à partir de ses observations. Les
coïncidences de valeurs sont accompagnées d’opérations,
de conditions de compatibilité et de critères de restitution.

La composition distingue le changement passif de coordonnées, qui
se télescope lorsque le système de coordonnées et la phase reviennent, du transport relatif
dont l’opérateur final est $H_{\mathrm{lazo}}\ne I$. Un échange
physique requiert de réaliser les opérations, leurs inverses et la
référence dans un même système, en incluant le contrôleur dans
le bilan. Ce critère conserve le contenu mathématique
d’une variation énergétique sans lui attribuer de production d’énergie
par le seul changement de lecture.

Les preuves de composition, d’itération, de restitution du coefficient
d’action, de réciprocité constitutive et de sensibilité soutiennent
ces caractérisations dans leurs domaines respectifs. La narration,
les définitions et les démonstrations forment ainsi un exposé continu :
la première identifie la signification des relations, les
deuxièmes fixent les objets et les troisièmes établissent les
opérations qui permettent de les composer et de les retrouver.
